\section{Probability estimates for the two certificate searches}
\label{cp:section}

The certificate search needs two facts about random formulas. First, it
needs computable lower bounds that approach the limiting density. Second,
it needs a nonnegative lower limit below the threshold and a negative
upper limit above it for normalized expected log partition functions,
with a suitable finite penalty in the latter case. These become pressure
signs when the pressure limit is established. The common
estimate controls the extra time for which a formula stays satisfiable
when one variable may be deleted, even if other deletions are already
allowed.

It is useful to begin with an auxiliary clause law. For every integer
$n\ge1$, clauses arrive in a Poisson process of rate one. Each of their three variable
indices is sampled independently and uniformly from $[n]$, and each sign
is sampled independently and fairly. All indices, signs, and the clock
are independent. Repetitions \emph{within} a clause are permitted; time
$an$ therefore gives a Poisson number of these clauses with mean $an$.
We call this the \emph{relaxed Poisson model}. Later in the section we
transfer its strict-side limits to the proper-clause model of the
introduction.

Fix throughout this section
\begin{equation}\label{cp:constants}
 B=20,\qquad p=\frac65,\qquad b=\frac1p=\frac56.
\end{equation}
For an integer $r$ with $0\le r\le n$, a formula is called $r$-deletion satisfiable if some set
of at most $r$ variables can be deleted so that the remaining formula is
satisfiable; deleting a variable discards every clause that touches it.
Let $T_r$ be the first time $r$-deletion satisfiability fails, with
$T_r=+\infty$ if it never fails, and put
\[
 F_r=T_r\wedge Bn.
\]
The formula at time $t$ includes all arrivals up to and including $t$.
Accordingly it survives at time $t$ exactly when $T_r>t$. For a fixed
variable $v$, a superscript $+v$ means that $v$ is deleted for free and up
to $r$ \emph{other} variables may be deleted. In particular,
$F_r^{+v}=T_r^{+v}\wedge Bn\ge F_r$, with the same cap $Bn$.
Allowing more deletions than there are other variables
has its literal meaning; it simply permits deleting all of them.
We suppress the dependence on $n$ in these times.  For this fixed size,
the overview's $T_n$ is $T_0$ here; both denote the uncapped first-unsatisfiable
time, whereas $F_0=T_0\wedge Bn$ is capped.  Write
\[
 f_n=\E F_0,\qquad P_n(t)=\Prob(T_0>t).
\]
Thus $f_n$ is the mean first-unsatisfiable time capped at $20n$.

\subsection{A uniform moment bound for one free deletion}

\begin{lemma}[Free-deletion moment]\label{cp:deletion-moment}
For all integers $n\ge1$ and $0\le r\le n$, and every $v\in[n]$,
\begin{equation}\label{cp:one-delay}
 \E\bigl[(F_r^{+v}-F_r)^p\bigr]\le C,\qquad C=2^{60}.
\end{equation}
Consequently, uniformly in $r$,
\begin{equation}\label{cp:concentration}
 \E\lvert F_r-\E F_r\rvert^p\le C_1n,
 \qquad C_1=2^{62}.
\end{equation}
\end{lemma}

The proof compares two ways to eliminate a set of surviving assignments.
Clauses containing the free variable once become two-literal tests after
that variable is fixed. A batch of relaxed three-literal tests can replace
each such test at controlled loss. This is the residual-set conditioning
principle used in Friedgut's half-cube lemma
\cite[Lemma~5.7]{Friedgut1999} and in Carenini's sequential clause
replacement \cite[Theorem~4.3 and Lemma~6.1]{Carenini}. The calculation
below is for independently sampled positions: their exact powers give
the estimate needed under an existing deletion allowance.

\begin{proof}
For $n=1$, Equation~\eqref{cp:one-delay} follows directly from
$0\le F_r^{+v}-F_r\le20$. Suppose henceforth that $n\ge2$.

\emph{Incident clauses.}
Split the Poisson process into the clauses avoiding $v$, called the
\emph{base}, and the clauses involving $v$. The base has rate
\[
 \lambda_n=\left(\frac{n-1}{n}\right)^3\ge\frac18,
\]
and its clauses have the relaxed distribution on the other $n-1$ variables.
It is independent of the incident streams.

At a deterministic time $t$, condition on the entire base up to $t$.
Let $U_t$ be its finite set of surviving pairs $(D,\sigma)$, where
$D\subseteq[n]\setminus\{v\}$ has size at most $r$ and $\sigma$ is an
assignment on $[n]\setminus(D\cup\{v\})$. When a literal is tested on such a
pair, a literal on $D$ is declared true. This convention exactly implements
clause deletion. The event $U_t\ne\varnothing$ is $\{T_r^{+v}>t\}$.
For a nonempty deterministic candidate set $U$, let $q(U)$ be the probability
that $\Poi(K)$ independent relaxed two-literal clauses on the other
variables eliminate every candidate, where $K=3B=60$.

For $0\le t<Bn$ and $U_t\ne\varnothing$, we claim
\begin{equation}\label{cp:incident-killing}
 \Prob(T_r\le t\mid\text{base up to }t)
 \le \bigl(q(U_t)+60/n\bigr)^2.
\end{equation}
To check this, fix a value for $v$. Only incident clauses in which every
occurrence of $v$ has the sign falsified by that value can obstruct an
extension of a candidate in $U_t$. Among these clauses, those containing
$v$ once project to independent relaxed two-literal clauses. Their number
is Poisson with mean at most $3t/n\le60$. The expected number containing
$v$ at least twice is at most $3t/n^2\le60/n$, by a union bound over the
three pairs of positions. The obstruction probability for this value is
therefore at most $q(U_t)+60/n$. The obstructing groups for the two values
of $v$ are disjoint Poisson thinnings, hence independent of one another
and of the base. If either group leaves a candidate, restoring $v$ with
that value gives an $r$-deletion satisfying assignment. Both groups must
therefore obstruct, which proves Equation~\eqref{cp:incident-killing}.

\emph{Replacing shorter tests.}
In any list of
independent tests, replacing one relaxed two-literal clause by $d\ge1$
independent relaxed three-literal clauses reduces the probability of
eliminating all candidates by at most $d^{-2}$. Indeed, condition on all
other tests, leaving a candidate set $U'$. There is nothing to prove if it
is empty. Otherwise let $x$ be the proportion of signed literals false on
every member of $U'$. Each variable contributes at most one such sign, so
$0\le x\le1/2$, including under the deleted-literal convention. The
two-literal clause alone eliminates $U'$ with probability $x^2$. The $d$
three-literal clauses do so with probability at least
$1-(1-x^3)^d\ge1-e^{-dx^3}$. If $dx^3>1$, this lower bound exceeds
$1/4\ge x^2$. If $dx^3\le1$, it is at least $dx^3/2$, and
\[
 x^2-\frac d2x^3\le\frac{16}{27d^2}<\frac1{d^2}.
\]
This proves the comparison. Applying it successively is valid because,
at each replacement, one can condition on all the other independent tests.

\emph{How long a residual set can survive.}
The replacement estimate converts a large value of $q(U_t)$ into a
chance that the future base eliminates $U_t$. Partition the possible
values $q(U_t)\ge1/n$ into disjoint dyadic bins
$q_0\le q(U_t)\le2q_0$, with $q_0=2^{-j}$ and $j\ge1$, using any
fixed endpoint convention. Every bin that is used has $q_0\ge1/(2n)$.
For its index $j$, set
\begin{equation}\label{cp:bin-parameters}
 L_j=j+122,\qquad
 d_j=\left\lceil\sqrt{4L_j/q_0}\right\rceil,\qquad
 \ell_j=32L_jd_j.
\end{equation}
Since $\E2^{\Poi(60)}=e^{60}<2^{120}$,
$\Prob(\Poi(60)>L_j)\le q_0/4$. If $U$ belongs to the bin, then
$L_j$ two-literal tests eliminate it with probability at least
$q(U)-q_0/4$. Replacing each test by $d_j$ three-literal tests loses at
most $L_j/d_j^2\le q_0/4$. Thus $L_jd_j$ three-literal tests eliminate
$U$ with probability at least $q_0/2$.

In a future interval of length $\ell_j$, the base has a Poisson number of
arrivals with mean at least $4L_jd_j$. If $M=L_jd_j\ge1$ and $X$ is
Poisson with mean $\mu\ge4M$, Chebyshev's inequality gives
\[
 \Prob(X<M)\le\frac{\mu}{(\mu-M)^2}
 \le\frac4{9M}<\frac12.
\]
Consequently this interval eliminates the fixed $U$ with probability at
least $q_0/4$. Write
\[
 S_{r,v}(t)=\Prob(T_r^{+v}>t),\qquad
 E_j(t)=\{U_t\ne\varnothing,\ q(U_t)\text{ belongs to bin }j\}.
\]
The independent future base then gives
\begin{equation}\label{cp:bin-time}
 \begin{split}
 S_{r,v}(t)-S_{r,v}(t+\ell_j)&\ge(q_0/4)\Prob(E_j(t)),\\
 \int_0^{Bn}\Prob(E_j(t))\,\dd t&\le4\ell_j/q_0.
 \end{split}
\end{equation}
For the second inequality, integrate the first: the integral of
$S_{r,v}(t)-S_{r,v}(t+\ell_j)$ over $[0,Bn]$ is at most $\ell_j$.
This argument does not require $T_r^{+v}$ to be finite. Conditional on the
base at time $t$ in $E_j(t)$, survival to $t+u$ is also bounded by
\begin{equation}\label{cp:bin-survival}
 (1-q_0/4)^{\lfloor u/\ell_j\rfloor}.
\end{equation}
For each complete future interval, failure to eliminate even the original
$U_t$ is necessary for survival; these interval events are independent.

\emph{Integrating the delay.}
For any $0\le F_r\le F_r^{+v}\le Bn$, direct integration gives the
pathwise identity
\begin{equation}\label{cp:delay-integral}
 (F_r^{+v}-F_r)^p
 =p(p-1)\int_{\substack{t,u\ge0\\t+u<Bn}}
 u^{p-2}\1_{\{T_r\le t,\ T_r^{+v}>t+u\}}\,\dd t\,\dd u.
\end{equation}
The boundaries have Lebesgue measure zero, so a hitting time equal to the
cap causes no difficulty. Conditional on the base up to $t$, incident
arrivals up to $t$ and future base arrivals are independent. The incident
killing estimate~\eqref{cp:incident-killing} and the future survival
estimate~\eqref{cp:bin-survival} therefore multiply.

Here is why the resulting sum over bins converges. On a bin of scale
$q_0$, incident killing has probability at most a constant times $q_0^2$.
The expected total time spent in that bin before $Bn$ is at most
$4\ell_j/q_0$, by
Equation~\eqref{cp:bin-time}. Integrating the geometric future-survival
bound against $u^{p-2}$ contributes at most a constant times
$(\ell_j/q_0)^{p-1}$. Thus the bin contribution is bounded by a constant times
$\ell_j^p q_0^{2-p}$. Since $\ell_j$ grows like a polynomial in $j$ times
$q_0^{-1/2}$, the power of $q_0$ is $2-3p/2=1/5>0$.

We now retain explicit constants for the later one-sided certificate.
On bin $j$,
\[
 (q(U_t)+60/n)^2\le(122q_0)^2<2^{14}q_0^2,
 \qquad
 \ell_j\le2^{20}(j+1)^{3/2}2^{j/2}.
\]
For the latter, $d_j\le3\sqrt{L_j/q_0}$ and
$L_j\le64(j+1)$ suffice. Moreover,
\[
 \begin{split}
 &\int_0^\infty u^{p-2}
     (1-q_0/4)^{\lfloor u/\ell_j\rfloor}\,\dd u\\
 &\qquad\le
 e\,\Gamma(1/5)(4\ell_j/q_0)^{1/5}
 <2^7(\ell_j/q_0)^{p-1}.
 \end{split}
\]
Indeed, the geometric term is at most
$e\exp(-q_0u/(4\ell_j))$, and splitting the defining Gamma integral
at $1$ gives $\Gamma(1/5)<6$. Using
Equation~\eqref{cp:bin-time} and $p(p-1)<1$, the contribution of bin $j$
to the expectation in Equation~\eqref{cp:delay-integral} is at most
\[
 2^{47}(j+1)^2\,2^{-j/5}.
\]
Grouping $j=5k+1,\ldots,5k+5$ gives
\[
 \sum_{j\ge1}(j+1)^2\,2^{-j/5}
 \le5\sum_{k\ge0}(5k+6)^2\,2^{-k}
 =1710<2^{11}.
\]
The total bin contribution is less than $2^{58}$.
When $q(U_t)<1/n$, Equation~\eqref{cp:incident-killing} is at most
$61^2/n^2$. Integrating over the entire truncated region in
Equation~\eqref{cp:delay-integral} bounds the remaining contribution by
\[
 61^2 B^p n^{p-2}<2^{18}.
\]
This proves Equation~\eqref{cp:one-delay} with $C=2^{60}$, without assuming
that the free-deletion problem eventually fails. In particular, if
$r\ge n-1$, a candidate can delete all other variables, so $q(U_t)=0$
and only the small-$q$ estimate is used.

\emph{From deletion to concentration.}
Partition all arrivals up to $Bn$ into $n$ independent
buckets according to the smallest variable index in a clause. Resample
bucket $i$ independently, and denote the resulting capped time by $F'_r$.
Deleting $i$ removes that entire bucket, so $F_r^{+i}$ is unchanged by
the resampling and dominates both $F_r$ and $F'_r$. Hence
\[
 \E|F_r-F'_r|^p
 \le\E(F_r^{+i}-F_r)^p+\E(F_r^{+i}-F'_r)^p
 \le2C.
\]
Reveal the buckets in their index order and let $D_i$ be the successive
differences of the Doob martingale for $F_r$. Averaging over the independent
replacement and all unrevealed buckets identifies $D_i$ as the conditional
expectation of $F_r-F'_r$ given the first $i$ buckets. Conditional Jensen
therefore gives $\E|D_i|^p\le2C$.

For real $x,y$ and $1<p<2$,
\[
 |x+y|^p\le |x|^p
 +p\operatorname{sgn}(x)|x|^{p-1}y+2|y|^p.
\]
Here $\operatorname{sgn}(0)=0$. The signed $(p-1)$-power is
H\"older continuous with constant
$2^{2-p}$, and integration of the derivative gives a remainder at most
$2^{2-p}|y|^p\le2|y|^p$. Apply this inequality successively to the
centered martingale sum. Each linear term has expectation zero, giving
$\E|F_r-\E F_r|^p\le2\sum_i\E|D_i|^p\le4Cn$.
This is Equation~\eqref{cp:concentration}.
\end{proof}

\subsection{Unequal blocks and a summable merging error}

The moment bound places the capped time close to its mean on a scale
$n^{5/6}$. To make the normalized means converge, we compare two
independent blocks with their union. The times assigned to the two
blocks may be arbitrary, so their densities need not agree.

\begin{lemma}[Weighted satisfiability interpolation]
\label{cp:interpolation}
For positive integers $n_1,n_2$ and real $t_1,t_2\ge0$,
\begin{equation}\label{cp:product-survival}
 P_{n_1+n_2}(t_1+t_2)
 \ge P_{n_1}(t_1)P_{n_2}(t_2).
\end{equation}
\end{lemma}

\begin{proof}
\emph{Mixing the blocks.}
Put $n=n_1+n_2$, $t=t_1+t_2$. The case $t=0$ is immediate, so set
$\lambda_i=t_i/t$. A \emph{local} clause first chooses block $i$ with
probability $\lambda_i$ and then chooses its three indices uniformly in
that block. A \emph{weighted global} clause chooses its block independently
at each of the three positions with these same probabilities, followed by
a uniform index in the chosen block. Interpolate between independent
Poisson counts of local and weighted global clauses, with respective means
$(1-s)t$ and $st$, for $0\le s\le1$.

The derivative of the expected satisfiability indicator is $t$ times the
expected difference of its add-one increments for the two clause types.
This follows by differentiating the Poisson series for a bounded function;
the differentiated series is absolutely convergent. For a satisfiable
base formula, let $x_i$ be the proportion of signed literals in block $i$
that are false in every satisfying assignment. A single clause eliminates
all satisfying assignments exactly when each of its literals is false
in every such assignment. The local and weighted global killing
probabilities are therefore
\[
 \sum_{i=1}^2\lambda_i x_i^3,
 \qquad \left(\sum_{i=1}^2\lambda_i x_i\right)^3,
\]
respectively. Convexity of $x^3$ on $[0,\infty)$ makes the derivative
nonnegative. An unsatisfiable base contributes zero. At the local endpoint
the two blocks are independent, so its satisfiability probability is the
right-hand side of Equation~\eqref{cp:product-survival}.

\emph{Balancing the variable weights.}
The mixed endpoint still has weight $\lambda_i/n_i$ on each variable
of block $i$. It remains to show that making these weights uniform cannot
decrease the probability at the weighted global endpoint. We prove this
for every fixed clause count. Choose two variables with weights $w_1,w_2$
and replace both weights by $(w_1+w_2)/2$. If their sum is zero, there is
nothing to change. Otherwise condition on the set of occurrence positions
that use one of these two variables, and on all choices and signs in the
other positions. These conditioned data have the same distribution before
and after averaging. Let $k$ be the number of selected positions. Treating
their literal truth values as independent placeholders defines a set
$S\subseteq\{0,1\}^k$: a vector belongs to $S$ if the variables outside
the selected pair can be assigned so that the formula is satisfied.
This set does not depend on the identities or signs still to be sampled
in the selected positions.

Let $y$ be the vector of literal truth values when the two selected
variables are both assigned zero, and let $\xi$ indicate positions occupied
by the second variable. Fair independent signs make $y$ uniform on the
bit cube, independently of $\xi$, whose coordinates are independent Bernoulli
variables with parameter $\theta=w_2/(w_1+w_2)$. The four assignments to
the selected pair yield exactly the truth vectors
\[
 y,\quad\bar y,\quad y\oplus\xi,\quad\overline{y\oplus\xi}.
\]
Here a bar denotes coordinatewise complement. Let
$A=S\cup\{\bar z:z\in S\}$ and $g=\1_A$. The conditional probability
of satisfaction is consequently
\[
 2\E g(y)-\E[g(y)g(y\oplus\xi)].
\]
Expand $g$ in the orthonormal parity characters
$\chi_I(y)=(-1)^{\sum_{j\in I}y_j}$ of the uniform bit cube. Independence
of the selector bits gives
\[
 \E[g(y)g(y\oplus\xi)]
 =\sum_{I\subseteq[k]}\widehat g(I)^2(1-2\theta)^{|I|}.
\]
Complement invariance of $g$ makes every odd-degree coefficient vanish.
Every remaining nonconstant term is nonnegative and is zero at
$\theta=1/2$. Thus equalizing the selected weights maximizes the
conditional satisfiability probability.

Repeatedly average a largest and a smallest weight. The sum over all
weights of squared deviations from $1/n$ decreases by half the square of their
difference. This nonnegative sum converges, so the largest-smallest
difference tends to zero and all weights converge to $1/n$. For a fixed
clause count the satisfiability probability is a polynomial in the
weights, hence continuous. It cannot decrease in this averaging procedure.
Bounded convergence then permits averaging over the Poisson count. The
uniform endpoint is the left-hand side of
Equation~\eqref{cp:product-survival}, proving the result.
\end{proof}

\begin{proposition}[A limiting center and a one-sided error]
\label{cp:lower}
The sequence $f_n/n$ converges to a real number $\alpha$. For every integer $k\ge1$,
\begin{equation}\label{cp:certified-lower}
 \alpha\ge\frac{f_k}{k}-2^{67}k^{-1/6}.
\end{equation}
\end{proposition}

\begin{proof}
We first turn the survival comparison into an almost-superadditive
inequality for $f_n$. Then a dyadic construction makes its accumulated
error summable after division by the size.
Set $H=2^{60}$ and $h_k=Hk^b$. Equation~\eqref{cp:concentration} and
Markov's inequality give
\[
 \Prob(|F_0-f_k|\ge h_k)
 \le C_1/H^p=2^{-10}<1/10
\]
at size $k$. In Lemma~\ref{cp:interpolation}, choose
$t_i=(f_{n_i}-h_{n_i})_+$. If $t_i=0$, survival is certain. Otherwise
$t_i<Bn_i$, so survival has probability at least $9/10$ by this deviation
bound. With $n=n_1+n_2$ and $t=t_1+t_2$, Equation~\eqref{cp:product-survival}
gives $P_n(t)\ge81/100$. Also $t<Bn$. If $f_n<t-h_n$, the same deviation
bound at size $n$ would give $P_n(t)\le1/10$, a contradiction. Hence
\begin{equation}\label{cp:almost-superadditive}
 f_{n_1+n_2}\ge f_{n_1}+f_{n_2}-C_2(n_1+n_2)^b,
 \qquad C_2=2^{62},
\end{equation}
because $h_{n_1}+h_{n_2}+h_n\le H(2^{1-b}+1)n^b<4Hn^b$.

Fix $k\ge1$. Iterating Equation~\eqref{cp:almost-superadditive} on
equal-size blocks gives
\[
 \frac{f_{2^jk}}{2^jk}
 \ge\frac{f_k}{k}
 -C_2k^{b-1}\sum_{i=1}^j2^{-(1-b)i}.
\]
Let
\[
 D_b=\sum_{i\ge1}2^{-(1-b)i}
 =\frac1{2^{1/6}-1}<32.
\]
The last inequality follows from $(33/32)^6<2$.
Thus every dyadic multiple $s=2^jk$ satisfies
$f_s\ge s(f_k/k-C_2D_bk^{-1/6})$.
For $n\ge k$, use the binary expansion of $\lfloor n/k\rfloor$ to
partition $k\lfloor n/k\rfloor$ into distinct dyadic multiples of $k$.
Merge them in increasing order of size. When a block of size $2^jk$ is
added, the total merged size is less than $2^{j+1}k$; the sum of merging
errors in Equation~\eqref{cp:almost-superadditive} is bounded by
\[
 C_2\sum_{j=0}^{\lfloor\log_2(n/k)\rfloor}(2^{j+1}k)^b
 \le A_b C_2 n^b
\]
for a fixed finite constant $A_b$. Add the remainder, if nonzero, in one
more application of Equation~\eqref{cp:almost-superadditive}, using
$f_s\ge0$ for its contribution. Divide by $n$ and let $n$ tend to
infinity with $k$ fixed. Since $b<1$, the merging errors vanish and
\[
 \liminf_{n\to\infty}\frac{f_n}{n}
 \ge\frac{f_k}{k}-C_2D_bk^{-1/6}
 \ge\frac{f_k}{k}-2^{67}k^{-1/6}.
\]
The bounded sequence $f_n/n\in[0,B]$ has a limsup subsequence. Let $k$
tend to infinity along that subsequence in the last inequality. Its
liminf is at least its limsup, proving convergence and
Equation~\eqref{cp:certified-lower}.
\end{proof}

\subsection{The threshold in the proper-clause model}

\begin{theorem}[Limiting satisfiability threshold]
\label{cp:threshold}
The limit $\alpha$ in Proposition~\ref{cp:lower} lies in $[1/200,10]$.
For every fixed real $a\ge0$, the relaxed Poisson model satisfies
\[
 P_n(an)\longrightarrow
 \begin{cases}
 1,&0\le a<\alpha,\\
 0,&a>\alpha.
 \end{cases}
\]
The same limits hold for $p(n,\lfloor an\rfloor)$ in the proper-clause
model defined in the introduction. In particular its limiting density is
$\alpha_3=\alpha$.
\end{theorem}

\begin{proof}
For a fixed assignment, a relaxed clause is violated with probability
$1/8$, even if some variable indices repeat, because the signs are
independent. The expected number of satisfying assignments at time $10n$
is therefore
\[
 2^n\exp(-10n/8)\longrightarrow0.
\]
Consequently $P_n(10n)\to0$, and
$f_n\le10n+(B-10)nP_n(10n)$ shows that $\alpha\le10<B$.

For a lower bound set $c=1/100$ and $M=\lfloor cn\rfloor$. A union bound
shows that the probability some $\ell$ of the first $M$ clauses touch
fewer than $\ell$ distinct variables is at most
\begin{equation}\label{cp:hall-bound}
 \sum_{\ell=1}^{M}\binom M\ell\binom n\ell
       (\ell/n)^{3\ell}
 \le\sum_{\ell=1}^{M}(e^2c\ell/n)^\ell=o(1).
\end{equation}
To obtain the first inequality, pad any set of fewer than $\ell$ touched
variables to one of size $\ell$; all $3\ell$ index choices must lie in
that set. The second uses $\binom u\ell\le(eu/\ell)^\ell$.
For the last limit, the terms with $\ell\le\sqrt n$ sum to
$O(n^{-1/2})$, and the remaining terms are at most
$n(e^2c^2)^{\sqrt n}$ in total. On the complementary event, Hall's
marriage theorem~\cite[Theorem~1]{Hall1935} assigns every clause a distinct representative variable
appearing in it. Set each representative to satisfy one of its occurrences
in its assigned clause. These requirements are consistent because the
representatives are distinct, and they satisfy all clauses. This argument
permits repeated indices and signs within a clause. At time $cn/2$ the
Poisson clause count is at most $M$ with probability tending to one.
Thus $P_n(cn/2)\to1$. Since $f_n\ge(cn/2)P_n(cn/2)$, we obtain
$\alpha\ge c/2=1/200$.

If $a<\alpha$, convergence of $f_n/n$ and
Equation~\eqref{cp:concentration} imply $P_n(an)\to1$. Explicitly, for
$a>0$ and sufficiently large $n$, $f_n-an\ge(\alpha-a)n/2$, so the
failure probability is at most
$C_1n/((\alpha-a)n/2)^p$, which tends to zero. The case $a=0$ is
immediate. If $\alpha<a<B$, the analogous upper-tail estimate gives
$P_n(an)\to0$. For $a\ge B$, use monotonicity and any intermediate
density in $(\alpha,B)$.

We now take $n\ge3$ and return to distinct-variable clauses. Write $P_n^{\mathrm d}(t)$
for their rate-one Poisson satisfiability probability. Thinning out the
relaxed clauses having repeated variable indices leaves an independent
distinct-clause stream of rate
\[
 s_n=(1-1/n)(1-2/n).
\]
Removing clauses can only help satisfiability. Conversely, the discarded
stream is empty up to time $an$ with probability
$\exp(-an(1-s_n))$, independently of the retained stream. Thus
\begin{equation}\label{cp:thinning}
 e^{-an(1-s_n)}P_n^{\mathrm d}(s_n a n)
 \le P_n(an)\le P_n^{\mathrm d}(s_n a n).
\end{equation}
For fixed $a$, the exponential factor is at least $e^{-3a}$.
If $a<\alpha$, choose $a'\in(a,\alpha)$. For large $n$, $s_na'>a$,
and monotonicity with Equation~\eqref{cp:thinning} gives
$P_n^{\mathrm d}(an)\ge P_n(a'n)\to1$.
If $a>\alpha$, choose $a'\in(\alpha,a)$. Then
\[
 P_n^{\mathrm d}(an)
 \le P_n^{\mathrm d}(s_na'n)
 \le e^{3a'}P_n(a'n)\longrightarrow0.
\]

Finally let $p(n,m)$ denote the probability for exactly $m$ independently
sampled distinct-variable clauses, allowing repeated clauses.
It is nonincreasing in $m$. Below the threshold, compare
$m=\lfloor an\rfloor$ with a Poisson count of mean $a'n$ for
$a<a'<\alpha$. This count is at least $m$ with probability tending to
one, so $p(n,m)\ge P_n^{\mathrm d}(a'n)-o(1)\to1$.
Above the threshold use $\alpha<a'<a$; the Poisson count is at most $m$
with probability tending to one, giving
$p(n,m)\le P_n^{\mathrm d}(a'n)+o(1)\to0$.
Conditioning three ordered uniform indices to be distinct and then
forgetting their order produces exactly a uniform three-element subset
with independent fair signs. Clauses remain
independent and may repeat. This completes the transfer to the exact model.
\end{proof}

\subsection{Computing the lower certificates}

The preceding theorem identifies the limit with the proper-clause
threshold. We now make the one-sided bound in
Equation~\eqref{cp:certified-lower} into finite rational certificates.

\begin{proposition}[Effective finite-size integration]
\label{cp:finite-computation}
The real $f_n$ is computable uniformly in the integer $n$. There is an
effective enumeration of rational numbers strictly below $\alpha$ whose
supremum is $\alpha$. In particular, there is a computable rational
sequence $(\ell_n)_{n\ge1}$ with $\ell_n<\alpha$ and
$\ell_n\to\alpha$.
\end{proposition}

\begin{proof}
For integers $n\ge1$ and $m\ge0$, let $s_{n,m}$ be the relaxed-model
satisfiability probability with exactly $m$ clauses. This is rational and
can be computed by enumerating all $(8n^3)^m$ ordered signed clause lists
and all $2^n$ assignments.
The survival-integral identity and conditioning on the Poisson count give
\[
 f_n=\int_0^{Bn}P_n(t)\,\dd t
 =\sum_{m=0}^\infty s_{n,m}
       \int_0^{Bn}e^{-t}\frac{t^m}{m!}\,\dd t.
\]
The summands are nonnegative. Uniformly for $0\le t\le Bn$,
\[
 \Prob(\Poi(t)>M)\le e^{Bn}2^{-(M+1)}.
\]
Thus truncation at $M$ changes the integral by at most
$Bn e^{Bn}2^{-(M+1)}$, a bound that can be made smaller than any prescribed
positive rational error by a finite search using the integer upper bound
$e^{Bn}<3^{Bn}$. Each remaining integral equals
\[
 1-e^{-Bn}\sum_{j=0}^m\frac{(Bn)^j}{j!},
\]
and is computable with certified rational intervals. For completeness,
if $x=Bn$ and $J\ge2x$ is an integer, the partial sum
$S_J=\sum_{j=0}^Jx^j/j!$ satisfies
\[
 S_J\le e^x\le S_J+2\frac{x^{J+1}}{(J+1)!}.
\]
Indeed, each ratio after the first omitted term is at most $1/2$.
The rational upper bound for the tail tends to zero. Taking reciprocals
gives certified intervals for $e^{-x}$, and then for each displayed
integral. Together with the explicit Poisson tail, this proves uniform
computability of $f_n$.

Rational bisection on $[0,1]$ computes the sixth root of $1/n$, so
\[
 v_n=\frac{f_n}{n}-2^{67}n^{-1/6}
\]
is uniformly computable. By Proposition~\ref{cp:lower}, $v_n\le\alpha$,
and $v_n\to\alpha$. For each pair of positive integers $(n,j)$, compute
a rational $\widehat v_{n,j}$ within $2^{-j}$ of $v_n$, and output
\[
 \ell_{n,j}=\widehat v_{n,j}-2^{1-j}.
\]
Every output is strictly less than $v_n$, hence less than $\alpha$.
Enumerating all pairs $(n,j)$ is effective, and the supremum of the
outputs is $\alpha$. To obtain a single sequence, choose effectively
$j(n)$ with $3\,2^{-j(n)}\le1/n$ and put $\ell_n=\ell_{n,j(n)}$.
Then
\[
 v_n-\frac1n\le\ell_n<v_n\le\alpha.
\]
Since $v_n\to\alpha$, this sequence also tends to $\alpha$.
It therefore supplies the lower sequence required by
Lemma~\ref{cert:search}.
\end{proof}

\subsection{Robust unsatisfiability and finite-temperature witnesses}

The lower certificate family is now complete. For the upper search, we
need a negative pressure witness above the threshold. Exact
unsatisfiability alone does not supply one: the number of violated
clauses could still be sublinear. The deletion moment prevents this by
bounding how much extra survival time a linear deletion allowance buys.

\begin{proposition}[Linear deletion and violation robustness]
\label{cp:robust}
For all integers $n\ge1$ and $0\le r\le n$,
\begin{equation}\label{cp:robust-mean}
 \frac{\E F_r}{n}\le\frac{f_n}{n}
       +C_4(r/n)^{1/6},\qquad C_4=2^{53}.
\end{equation}
For every fixed real $a>\alpha$, there exists $\epsilon>0$ such that, with probability
tending to one, the relaxed formula at time $an$ cannot be made satisfiable
by deleting $\lfloor\epsilon n\rfloor$ variables. On the same event every
assignment violates more than $\lfloor\epsilon n\rfloor$ clauses.
\end{proposition}

\begin{proof}
The case $r=0$ is equality. For $r\ge1$, fix a realization.
The survival time with at most $r$ deletions is
the maximum of the survival times over finitely many fixed deletion sets.
Capping commutes with this finite maximum. Choose a maximizing set $D$,
padding it to size exactly $r$; padding cannot worsen survival. For every
$v\in D$, the free deletion of $v$ followed by at most $r-1$ other
deletions has capped survival time exactly $F_r$: the choice $D\setminus
\{v\}$ attains it, and every competing choice uses at most $r$ deletions
in total. Hence, for $1\le r\le n$,
\[
 r(F_r-F_{r-1})^p
 \le\sum_{v=1}^n(F_{r-1}^{+v}-F_{r-1})^p.
\]
This includes the case of survival up to the cap and the case $r=n$.
Take expectations and use Lemma~\ref{cp:deletion-moment}, followed by
the $L^p$ bound on the first moment. Summing over the deletion allowance,
\[
 \E F_r-f_n
 \le (Cn)^{1/p}\sum_{j=1}^r j^{-1/p}
 \le6C^{5/6}n^{5/6}r^{1/6}
 <2^{53}n^{5/6}r^{1/6},
\]
which proves Equation~\eqref{cp:robust-mean}.

Fix $a>\alpha$ and choose $a_0\in(\alpha,\min(a,B))$. Let
$\Delta=a_0-\alpha>0$ and choose $0<\epsilon<1$ so small that
$C_4\epsilon^{1/6}<\Delta/4$. For large $n$, convergence of $f_n/n$
and Equation~\eqref{cp:robust-mean}, with $r=\lfloor\epsilon n\rfloor$,
give $\E F_r\le(\alpha+\Delta/2)n$. Since $a_0<B$, survival at time
$a_0n$ implies $F_r>a_0n$. Equation~\eqref{cp:concentration} therefore
shows
\[
 \Prob(T_r>a_0n)
 \le\frac{C_1n}{(\Delta n/2)^p}\longrightarrow0.
\]
By monotonicity the same conclusion holds at time $an$. If an assignment
violated at most $r$ clauses, selecting one variable from each violated
clause and deleting those variables would discard every violation. The
remaining clauses would be satisfied by the restricted assignment, using
at most $r$ deletions. This contradiction proves the violation assertion.
\end{proof}

For a relaxed formula at time $an$, let $H_n(\sigma)$ be the number of
clauses violated by $\sigma$, and define
\[
 Z_n(a,\beta)=\sum_{\sigma\in\{0,1\}^n}
       e^{-\beta H_n(\sigma)},\qquad \beta>0.
\]

\begin{corollary}[Signs of the soft pressure]\label{cp:soft-signs}
For every fixed real $0\le a<\alpha$ and every fixed $\beta>0$,
\begin{equation}\label{cp:soft-below}
 \liminf_{n\to\infty}\frac1n\E\log Z_n(a,\beta)\ge0.
\end{equation}
For every fixed real $a>\alpha$, some positive integer $\beta$ satisfies
\begin{equation}\label{cp:soft-above}
 \limsup_{n\to\infty}\frac1n\E\log Z_n(a,\beta)<0.
\end{equation}
\end{corollary}

These statements concern the finite-size expected logarithms and do not
need existence of a pressure limit. Once that limit is proved, they give
the two signs used in the upper certificate search.

\begin{proof}
Put $V_n=\min_\sigma H_n(\sigma)$, and let $M_n\sim\Poi(an)$ be the
total clause count. Below the threshold, $\Prob(V_n>0)\to0$ and
$0\le V_n\le M_n$. Cauchy--Schwarz gives
\[
 \frac{\E V_n}{n}
 \le\bigl((a^2+a/n)\Prob(V_n>0)\bigr)^{1/2}\longrightarrow0.
\]
The bound $\log Z_n\ge-\beta V_n$ proves
Equation~\eqref{cp:soft-below}. Above the threshold,
Proposition~\ref{cp:robust} gives $\liminf_n\E V_n/n\ge\epsilon$
for some $\epsilon>0$. Since
$\log Z_n\le n\log2-\beta V_n$, any integer
$\beta>(\log2)/\epsilon$ proves Equation~\eqref{cp:soft-above}.
\end{proof}
